% =============================================================================
% Construyendo un cerebro digital para la investigación en libre competencia:
% una guía práctica con Graphify para el caso chileno
%
% VERSIÓN CORREGIDA — Mayo 2026
% Revisión integral de estructura, redacción y formato sobre el borrador original
% (main.tex). El resumen de cambios figura al final del archivo.
%
% Compilación:
%   pdflatex main_corregido.tex
%   bibtex main_corregido
%   pdflatex main_corregido.tex
%   pdflatex main_corregido.tex
% =============================================================================

\documentclass[12pt,a4paper]{article}

% ----- Codificación y lengua -----
\usepackage[utf8]{inputenc}
\usepackage[T1]{fontenc}
\usepackage[spanish,es-nodecimaldot]{babel}

% ----- Tipografía -----
\usepackage{lmodern}
\usepackage{microtype}

% ----- Geometría -----
\usepackage[top=3cm,bottom=3cm,left=3.5cm,right=3.5cm]{geometry}

% ----- Matemáticas -----
\usepackage{amsmath}

% ----- Gráficos y colores -----
\usepackage{graphicx}
\usepackage[dvipsnames,table]{xcolor}
\definecolor{darkblue}{RGB}{0,51,102}
\definecolor{lightgray}{RGB}{245,245,245}
\definecolor{tabhead}{RGB}{30,70,120}
\definecolor{accent}{RGB}{160,40,40}

% ----- Tablas -----
\usepackage{booktabs}
\usepackage{longtable}
\usepackage{array}
\usepackage{colortbl}
\usepackage{tabularx}

% ----- Notas tipo "pendiente" (placeholders elegantes) -----
% Habilitar/deshabilitar todos los placeholders con disable/colorinlistoftodos
\usepackage[textsize=footnotesize,backgroundcolor=yellow!18,
            linecolor=accent,bordercolor=accent!60]{todonotes}
\newcommand{\pend}[1]{\todo[inline,caption={Pendiente}]{\textbf{Pendiente:} #1}}

% ----- Recuadros -----
\usepackage{tcolorbox}
\tcbuselibrary{skins,breakable}
\newtcolorbox{infobox}[1][]{
  colback=lightgray, colframe=darkblue!60,
  fonttitle=\bfseries\small, title=#1,
  breakable, sharp corners
}
\newtcolorbox{keybox}[1][]{
  enhanced, colback=darkblue!4, colframe=darkblue,
  fonttitle=\bfseries, title=#1, coltitle=white,
  attach boxed title to top left={xshift=6mm,yshift=-3mm},
  boxed title style={colback=darkblue, sharp corners},
  breakable, sharp corners
}
\newtcolorbox{codebox}{
  colback=black!3, colframe=black!30,
  fontupper=\ttfamily\small, breakable, sharp corners
}

% ----- Listas y enumeraciones -----
\usepackage{enumitem}
\setlist[itemize]{noitemsep, topsep=4pt}
\setlist[enumerate]{noitemsep, topsep=4pt}

% ----- Encabezados y secciones -----
\usepackage{titlesec}
\titleformat{\section}{\large\bfseries\color{darkblue}}{}{0em}{}[\titlerule]
\titleformat{\subsection}{\normalsize\bfseries\color{darkblue!80}}{}{0em}{}
\titleformat{\subsubsection}{\normalsize\itshape}{}{0em}{}

% ----- Hyperref (siempre al final, antes de biblatex) -----
\usepackage[colorlinks=true,linkcolor=darkblue,citecolor=darkblue,
            urlcolor=darkblue,pdftitle={Cerebro Digital Libre Competencia Chile}]{hyperref}

% ----- Notas al pie -----
\usepackage{footmisc}

% ----- Referencias -----
\usepackage[backend=bibtex,style=footnote-dw,autocite=footnote]{biblatex}
\addbibresource{references.bib}

% ----- Metadatos -----
\title{%
  \textbf{Construyendo un cerebro digital\\
  para la investigación en libre competencia}\\[0.4em]
  \large\textit{Una guía práctica con Graphify para el caso chileno}
}
\author{%
  \pend{Nombre del autor y afiliación}\\[0.3em]
  \small \texttt{\pend{correo electrónico}}
}
\date{Mayo 2026\\[0.5em]
  \small Borrador de trabajo — comentarios bienvenidos}

% =============================================================================
\begin{document}
% =============================================================================

\maketitle
\thispagestyle{empty}

\begin{abstract}
\noindent
Un \emph{cerebro digital} es una infraestructura personal de conocimiento
construida sobre documentos que el investigador selecciona y confía. Mapea las
conexiones entre fuentes, expone patrones e inconsistencias, y produce respuestas
trazables al material subyacente. Karpathy describió la práctica en su forma
mínima; Schrepel formalizó la metodología para investigación
jurídica.\autocite{schrepel2026brain}

Este artículo documenta la adaptación de esa metodología al derecho de la libre
competencia chileno. Sobre un corpus de aproximadamente 520 documentos
—sentencias del TDLC y la Corte Suprema, resoluciones, acuerdos extrajudiciales,
instrucciones de carácter general, expedientes de recomendación normativa, y
normativa de base (DL~211 y guías FNE)— se construye un grafo de conocimiento
consultable y una wiki temática. La implementación usa Graphify\autocite{shamsi2025graphify}
para el grafo, MarkItDown\autocite{microsoft2024markitdown} para la conversión, y
Claude Code como entorno de consulta. El archivo \texttt{CLAUDE.md} concentra el
diseño del esquema: jerarquía de autoridades, ponderación por centralidad de nodo,
anclas conceptuales y extracción a nivel de proposición.

El artículo distingue dos instrumentos complementarios construibles sobre el
mismo corpus. El \textbf{grafo de conocimiento} sirve a la investigación
estructural: identificar \emph{god nodes}, caminos ausentes, contradicciones
TDLC/CS y silencios doctrinales. El \textbf{sistema RAG vectorial} sirve a la
práctica diaria del abogado: localizar precedentes concretos y citar
literalmente. Ambos responden a preguntas distintas; este trabajo se centra en el
primero y trata el segundo como extensión.

El pipeline es replicable, de código abierto y neutro respecto del tamaño del
corpus. El repositorio quedará disponible públicamente.
\end{abstract}

\vspace{1em}
\noindent\textbf{Palabras clave:} grafo de conocimiento; libre competencia;
jurisprudencia chilena; TDLC; FNE; Graphify; \emph{god nodes}; RAG vectorial;
DL~211.

\newpage
\tableofcontents
\newpage

% =============================================================================
\section{Introducción}
% =============================================================================

Las instituciones de competencia acumulan documentos durante décadas. Sentencias,
resoluciones, investigaciones, guías, estudios sectoriales e informes periciales
se apilan en registros públicos. Cada pieza nació en un contexto particular y,
con frecuencia, en equipos sin visibilidad sobre lo escrito antes. El cuerpo
resultante es rico en sustancia pero opaco en estructura: las conexiones entre
documentos son difíciles de captar, la evolución de las posiciones cuesta
trazarla, y las inconsistencias suelen pasar inadvertidas.

El investigador enfrenta el mismo problema. Quien trabaja en la frontera del
campo ha leído cientos de sentencias y citado muchas. Leer y citar, sin embargo,
no equivale a retener un argumento. Al momento de escribir, las conexiones que
parecían nítidas durante la lectura se han disuelto; cuanto más antiguo el
documento, mayor la pérdida. Una sentencia leída hace diez años sobrevive como
conclusión, rara vez como razonamiento.\autocite{schrepel2026brain} Un grafo de
conocimiento aborda esto directamente: externaliza la estructura del corpus,
hace visibles y consultables las conexiones, y libera al escritor de la memoria.

Este artículo describe un pipeline que aplica esa idea al derecho chileno de la
libre competencia. Toma una carpeta de PDFs —sentencias del TDLC y la Corte
Suprema, resoluciones, acuerdos extrajudiciales, instrucciones de carácter
general, expedientes de recomendación normativa y normativa de base— y produce
tres salidas: (i)~una versión Markdown estandarizada de cada documento;
(ii)~un grafo de conocimiento con conceptos, relaciones y patrones de
razonamiento extraídos del corpus; y (iii)~una wiki organizada por comunidad
temática, con hipervínculos internos y un índice maestro. El sistema permite
consultar el corpus en lenguaje natural: cada respuesta cita la fuente.

\subsection{Dos instrumentos para dos preguntas distintas}

Conviene fijar una distinción que la literatura sobre IA aplicada al derecho
suele difuminar: la diferencia entre un \textbf{grafo de conocimiento} y un
\textbf{sistema RAG vectorial}. Ambos pueden construirse sobre el mismo corpus.
No son intercambiables. Responden a preguntas distintas y entregan resultados de
naturaleza distinta.

\begin{keybox}[Graphify para investigación vs.\ RAG vectorial para el día a día]
\textbf{Grafo (Graphify).} Extrae nodos y aristas del corpus completo, los
clusteriza en comunidades temáticas y permite preguntar por la \emph{estructura}
del campo: qué documentos son centrales, qué clusters están aislados, qué
caminos esperados no existen, dónde hay contradicciones entre niveles de
autoridad. Su valor es \emph{generativo}: produce hipótesis y agenda de
investigación, no fragmentos textuales.

\medskip
\textbf{RAG vectorial.} Indexa fragmentos de texto en un espacio de \emph{embeddings},
recupera los más similares semánticamente a la consulta, y los pasa a un modelo
de lenguaje para que sintetice una respuesta con cita literal. Su valor es
\emph{recuperativo}: encuentra el considerando exacto, el párrafo que define un
término, la línea de jurisprudencia que respalda un argumento. Es el instrumento
natural para preparar escritos.

\medskip
\textbf{Regla práctica.} Si la pregunta es ``¿qué dice el TDLC sobre X?'', use
RAG. Si la pregunta es ``¿qué \emph{no} dice el TDLC sobre X?'', use el grafo.
La Tabla~\ref{tab:grafo-vs-rag} desarrolla la comparación. Este artículo se
centra en el grafo; el sistema RAG aparece como extensión natural en
la~Sección~\ref{sec:discusion}.
\end{keybox}

\begin{table}[h!]
\centering
\caption{Grafo de conocimiento vs.\ sistema RAG vectorial}
\label{tab:grafo-vs-rag}
\small
\begin{tabularx}{\linewidth}{>{\bfseries}p{3.3cm} X X}
\toprule
\rowcolor{tabhead}\color{white}\textbf{Dimensión} &
  \color{white}\textbf{Grafo (Graphify)} &
  \color{white}\textbf{RAG vectorial (ChromaDB + API)} \\
\midrule
Tipo de pregunta &
  Estructural: god nodes, vacíos, contradicciones, evolución doctrinal &
  Textual: ¿qué dice el Considerando 14? ¿qué fallos mencionan ``margin squeeze''? \\
\addlinespace
Formato de salida &
  Listas de nodos, caminos, comunidades, ausencias &
  Fragmentos con cita exacta al documento fuente \\
\addlinespace
Caso de uso &
  Investigación: generar hipótesis, detectar silencios &
  Litigio y asesoría: redactar escritos, buscar precedentes \\
\addlinespace
Riesgo de alucinación &
  Bajo: el grafo acota el espacio de respuestas &
  Moderado: la cita debe verificarse contra el documento \\
\addlinespace
Costo computacional &
  Alto en construcción; bajo en consulta &
  Moderado en embeddings; bajo en consulta \\
\addlinespace
Tiempo de consulta &
  Inmediato una vez construido el grafo &
  Inmediato sobre el índice vectorial \\
\addlinespace
Foco de este artículo &
  Sí &
  Extensión futura (Sección~\ref{sec:discusion}) \\
\bottomrule
\end{tabularx}
\end{table}

\subsection{Por qué el caso chileno}

El sistema chileno tiene una dualidad institucional que lo vuelve un caso
exigente para este tipo de herramientas. El TDLC es un tribunal especializado
con ministros economistas; la Corte Suprema, al conocer recursos contra el TDLC,
aplica una interpretación generalista. La tensión es persistente: doctrinas
confirmadas por el especialista pueden ser modificadas o revocadas por el
generalista. Un grafo bien construido puede rastrear ese cruce de
forma sistemática —algo que un buscador de texto no captura, porque la pregunta
no es por ocurrencias de un término sino por la relación entre dos
pronunciamientos.

La FNE añade una tercera perspectiva: investigación administrativa que precede
al litigio. El corpus contiene 2.641 documentos de enforcement aún pendientes de
incorporación al grafo, lo que convierte este trabajo en una primera etapa de
un proyecto más amplio.

% =============================================================================
\section{Metodología del pipeline}
% =============================================================================

El pipeline tiene cinco pasos y un régimen de mantenimiento. Cada paso alimenta
al siguiente; el proceso puede interrumpirse y reanudarse sin pérdida de
progreso.\autocite{schrepel2026brain} El Anexo~\ref{anx:claude-md} reproduce la
versión anonimizada del archivo \texttt{CLAUDE.md} que gobierna la extracción.

\begin{figure}[h!]
  \centering
  \fbox{\parbox{0.88\linewidth}{\centering\vspace{0.8em}
    \textbf{Figura 1 (pendiente):} Diagrama de flujo del pipeline\\[0.4em]
    \texttt{Paso 1 (Recolectar) $\rightarrow$ Paso 2 (Convertir) $\rightarrow$
            Paso 3 (Grafo) $\rightarrow$}\\
    \texttt{Paso 4 (Wiki) $\rightarrow$ Paso 5 (Consultar)
            $\leftarrow$ Mantenimiento}\\
    \vspace{0.5em}
    Reemplazar por \texttt{figures/pipeline\_diagram.pdf}
  \vspace{0.4em}}}
  \caption{Pipeline de construcción del cerebro digital}
  \label{fig:pipeline}
\end{figure}

\subsection{Paso 1 — Recolectar}

Descargar los PDFs desde los repositorios públicos de cada institución:

\begin{itemize}
  \item \textbf{TDLC} (\texttt{tdlc.cl}): sentencias, resoluciones, acuerdos
    extrajudiciales (AE), instrucciones de carácter general (ICG), expedientes
    de recomendación normativa (ERN).
  \item \textbf{Corte Suprema} (\texttt{pjud.cl}): sentencias y resoluciones que
    conocen recursos contra el TDLC.
  \item \textbf{FNE} (\texttt{fne.cl}): investigaciones, resoluciones de archivo,
    guías y reglamentos.
\end{itemize}

Para los Tracks A (jurisprudencia) y C (normativa de base), la descarga puede
automatizarse con scripts en Python análogos al desarrollado por Schrepel para
las decisiones de la Comisión Europea.\autocite{schrepel2026brain} El Track~B
(2.641 documentos de enforcement) requiere un \emph{scraper} específico,
disponible en la subcarpeta \texttt{code/} del repositorio.

Los archivos se organizan en subcarpetas temáticas bajo \texttt{raw/}, espejo
exacto de la estructura que reproducirá \texttt{raw-md/} tras la conversión.

\subsection{Paso 2 — Convertir con MarkItDown}

Los PDFs no pueden ir directos al grafo. El asistente leería ruido de formato
—saltos de página, columnas, encabezados repetidos— en lugar de contenido.
MarkItDown, de Microsoft, convierte el documento a Markdown preservando la
estructura semántica (títulos, listas, tablas) y descartando el formato
visual.\autocite{microsoft2024markitdown}

\begin{codebox}
pip3 install 'markitdown[pdf]'
python3 code/batch_convert.py
\end{codebox}

El script es idempotente (\texttt{SKIP\_DONE = True}): al re-ejecutarlo, solo
procesa archivos nuevos o modificados. Los Markdown resultantes se depositan en
\texttt{raw-md/} replicando la estructura de \texttt{raw/}.

Para PDFs escaneados existe una variante con OCR
(\texttt{code/batch\_convert\_ocr.py}). El rendimiento sobre sentencias antiguas
del TDLC —típicamente anteriores a 2010— es desigual; conviene revisar
manualmente los archivos registrados en \texttt{conversion\_errors.log}.

\subsection{Paso 3 — Construir el grafo con Graphify}

Graphify se ejecuta dentro de Claude Code y lee la carpeta de
Markdown.\autocite{shamsi2025graphify} Opera en tres pasadas. La primera, AST
determinista, extrae estructura sin IA. La segunda despacha agentes en paralelo
—hasta 31 simultáneos— que extraen conceptos y relaciones de los documentos
asignados y devuelven JSON estructurado. La tercera fusiona los \emph{payloads}
en un grafo NetworkX, lo clusteriza con detección de comunidades Leiden, y lo
exporta en tres formatos: HTML interactivo, JSON consultable y un informe de
auditoría en Markdown.

\begin{codebox}
# Dentro de Claude Code, apuntando a la carpeta de Markdown:
/graphify raw-md/
# Cuando se consulta el alcance: all
# En sesiones subsiguientes (retoma desde el caché SHA-256):
/graphify raw-md/ --update
\end{codebox}

La caché SHA-256 evita reprocesar lo ya hecho: en sesiones posteriores, solo los
archivos nuevos o modificados pasan a los agentes. Para un corpus de
520~documentos como el Track~A+C, la extracción completa toma entre 15 y
30~minutos, repartidos en una o dos sesiones.

\begin{infobox}[Lecciones de operación: concurrencia y deduplicación]
La experiencia sobre este corpus sugiere que \texttt{max\_concurrency=4} es un
valor conservador que evita errores de límite de tasa de la API. Con cupos más
altos puede subirse a 8 o 16. Tras cada ejecución conviene una pasada de
deduplicación: Graphify ocasionalmente crea nodos con etiquetas próximas para el
mismo concepto (por ejemplo, ``mercado relevante'' frente a ``definición de
mercado relevante''). La consolidación puede hacerse manualmente o instruirse
mediante un \emph{prompt} de limpieza a Claude Code.
\end{infobox}

\subsection{Paso 4 — Generar la wiki}

\begin{codebox}
/graphify raw-md/ --wiki
\end{codebox}

Cada comunidad temática del grafo se convierte en un artículo de wiki. La
ponderación por centralidad de nodo, codificada en \texttt{CLAUDE.md}, asegura
que el resumen de cada artículo se construya primero desde los \emph{god nodes}
del cluster, no desde el promedio del corpus. Cada artículo sigue la estructura
fija descrita en la Sección~\ref{sec:esquema}: (1)~Resumen, (2)~Jurisprudencia
Track~A, (3)~Enforcement FNE Track~B, (4)~Casos vinculados, (5)~Lo que las
fuentes no abordan.

Los artículos se depositan en \texttt{graphify-out/wiki/} y pueden navegarse en
Obsidian, que renderiza los hipervínculos entre páginas y ofrece una vista de
grafo del corpus wiki.

\subsection{Paso 5 — Consultar}

\begin{codebox}
graphify claude install
\end{codebox}

Tras la instalación, Claude Code navega el grafo antes de responder cualquier
pregunta sobre el corpus. Cada afirmación es trazable a un documento.
El protocolo de consulta de seis pasos para generar hipótesis se describe
en la Sección~\ref{sec:protocolo}, y el Anexo~\ref{anx:consultas} contiene
los \emph{prompts} de ejemplo.

\subsection{Mantenimiento}

Para incorporar documentos nuevos basta depositar el PDF en la subcarpeta
temática correspondiente y reejecutar las tres operaciones:

\begin{codebox}
markitdown new_doc.pdf -o raw-md/sentencias-tdlc/new_doc.md
/graphify raw-md/ --update
/graphify raw-md/ --wiki
\end{codebox}

Solo se reprocesan los archivos nuevos o modificados. El grafo existente se
conserva; los nodos y aristas adicionales se integran y las comunidades se
re-clusteran.

% =============================================================================
\section{Implementación en el corpus chileno}
\label{sec:esquema}
% =============================================================================

Esta sección describe el corpus, las particularidades institucionales que el
esquema debe codificar, y las decisiones de diseño del \texttt{CLAUDE.md} que
distinguen un sistema de recuperación de información de un instrumento de
generación de hipótesis.

\subsection{Composición del corpus}

El corpus opera en tres tracks con lógicas distintas. La Tabla~\ref{tab:track-a}
detalla la composición del Track~A (jurisprudencia) y el Track~C (normativa de
base) ya procesados. La Tabla~\ref{tab:track-b} cubre el Track~B (enforcement
FNE) pendiente de incorporación.

\begin{table}[h!]
\centering
\caption{Tracks A y C — corpus procesado en el grafo (520 documentos)}
\label{tab:track-a}
\small
\begin{tabularx}{\linewidth}{l l >{\raggedleft\arraybackslash}X}
\toprule
\rowcolor{tabhead}\color{white}\textbf{Subcarpeta en \texttt{raw-md/}} &
  \color{white}\textbf{Tipo} &
  \color{white}\textbf{N°} \\
\midrule
\multicolumn{3}{l}{\itshape Track A — Jurisprudencia} \\
\texttt{sentencias-tdlc-md/} & Sentencias TDLC & 212 \\
\texttt{sentencias-cs-md/} & Sentencias Corte Suprema & 128 \\
\texttt{resoluciones-tdlc-md/} & Resoluciones TDLC & 89 \\
\texttt{resoluciones-cs-md/} & Resoluciones Corte Suprema & 13 \\
\texttt{ae-tdlc-md/} & Acuerdos Extrajudiciales (FNE–TDLC) & 38 \\
\texttt{ern-tdlc-md/} & Expedientes de Recomendación Normativa & 22 \\
\texttt{icg-tdlc-md/} & Instrucciones de Carácter General & 9 \\
\midrule
\multicolumn{3}{l}{\itshape Track C — Normativa de base} \\
\texttt{normativa-fusiones-fne-md/} & Guías, reglamentos y formularios FNE & 8 \\
\texttt{normativa-general-md/} & DL~211 refundido & 1 \\
\midrule
\multicolumn{2}{r}{\textbf{Total procesado}} & \textbf{520} \\
\bottomrule
\end{tabularx}
\end{table}

\begin{table}[h!]
\centering
\caption{Track B — Enforcement FNE (pendiente de incorporación al grafo)}
\label{tab:track-b}
\small
\begin{tabularx}{\linewidth}{l l >{\raggedleft\arraybackslash}X}
\toprule
\rowcolor{tabhead}\color{white}\textbf{Subcarpeta en \texttt{raw/}} &
  \color{white}\textbf{Tipo} &
  \color{white}\textbf{N° aprox.} \\
\midrule
\texttt{concentraciones-o-integraciones/} & Decisiones FNE & 1.341 \\
\texttt{abusos-de-posicion-de-dominio/} & Investigaciones FNE & 561 \\
\texttt{acuerdos-colusorios/} & Investigaciones FNE & 193 \\
\texttt{actos-anticompetitivos-de-la-autoridad/} & Investigaciones FNE & 146 \\
\texttt{incumplimiento-de-sentencias-e-instrucciones/} & Investigaciones FNE & 146 \\
\texttt{restricciones-verticales/} & Investigaciones FNE & 36 \\
\texttt{competencia-desleal/} & Investigaciones FNE & 52 \\
\texttt{sin-categoria/} & Varios & 73 \\
\midrule
\multicolumn{2}{r}{\textbf{Total pendiente}} & \textbf{2.641} \\
\bottomrule
\end{tabularx}
\end{table}

Cada track cumple un papel distinto. El Track~A es la base del análisis
jurisprudencial y la preparación de litigios: la jerarquía de autoridad aplica
plenamente. El Track~B muestra el estándar práctico con que la FNE evalúa
problemas de competencia —qué investiga, qué metodologías aplica, cuándo
archiva, cuándo requiere— y es indispensable para anticipar su comportamiento
como contraparte o requirente. El Track~C funciona como ancla normativa: los
nodos del Track~A deben conectarse a los artículos del DL~211 que invocan.

\subsection{La tensión TDLC / Corte Suprema como objeto de extracción}

La dualidad institucional chilena es estructuralmente relevante para el diseño
del esquema. El TDLC opera con un enfoque interdisciplinario, con ministros
economistas; la Corte Suprema aplica una interpretación generalista. Cuando la
CS revoca al TDLC por razones probatorias —exige prueba directa donde el TDLC
admitió prueba indirecta— la divergencia no es un episodio: es un patrón
doctrinal con consecuencias estratégicas para el litigio. El grafo puede
identificarlo de forma sistemática.

\begin{infobox}[Casos vinculados]
Una investigación FNE que deriva en requerimiento ante el TDLC y luego en
pronunciamiento de la CS produce tres documentos que son perspectivas del mismo
conflicto. El esquema instruye al sistema a identificarlos como caso vinculado y
tratarlos integradamente. La señal de vinculación son las referencias textuales
en las sentencias y resoluciones al rol de la investigación FNE de origen.
\end{infobox}

\subsection{El archivo CLAUDE.md como diseño del esquema}

\texttt{CLAUDE.md} es el único punto de intervención del investigador en el
proceso de extracción. Lo que contiene se incorpora al contexto de trabajo de la
IA cuando construye el grafo. Usado con cuidado, determina si el sistema
resultante recupera información o genera hipótesis.\autocite{schrepel2026brain}
La versión completa, anonimizada, está en el Anexo~\ref{anx:claude-md}. Cuatro
decisiones de diseño cambian cualitativamente lo que produce.

\subsubsection*{(a) Ponderación de autoridad}

Los corpora jurídicos contienen documentos con autoridad desigual. El esquema
debe codificar esa jerarquía explícitamente. Para el sistema chileno:

\begin{enumerate}[label=\arabic*.]
  \item Sentencias y resoluciones de la \textbf{Corte Suprema}: vinculantes;
    pueden confirmar, revocar o modificar al TDLC.
  \item Sentencias y resoluciones del \textbf{TDLC}: tribunal especializado;
    supeditado a la CS.
  \item[\phantom{0}2b.] \textbf{Instrucciones de Carácter General (ICG)} del
    TDLC, bajo Art.~18 N°3 del DL~211: vinculantes para todos los agentes del
    mercado.
  \item[\phantom{0}2c.] \textbf{Acuerdos Extrajudiciales (AE)} aprobados por el
    TDLC: vinculantes para las partes; señal del estándar de enforcement.
  \item[\phantom{0}2d.] \textbf{Expedientes de Recomendación Normativa (ERN)}:
    no vinculantes; revelan el análisis del tribunal sobre mercados regulados.
  \item Decisiones y dictámenes de la \textbf{FNE}.
  \item Guías y reglamentos FNE: peso auxiliar en litigios.
  \item Estudios sectoriales de la FNE.
  \item Doctrina académica revisada por pares.
  \item[0.] Notas personales del investigador: nivel cero; sin autoridad jurídica.
\end{enumerate}

A la jerarquía se suma una \textbf{regla de tensión TDLC/CS obligatoria}: toda
respuesta que cite una sentencia del TDLC debe indicar si la CS se pronunció
—confirmando, modificando, revocando, o sin haber sido recurrida—. La regla
aplica incluso cuando la tensión es menor.

\subsubsection*{(b) Los silencios como objeto consultable}

Los esquemas estándar instruyen al sistema a resumir lo que las fuentes dicen.
Un esquema de investigación debe instruirlo a identificar lo que las fuentes
\emph{no} abordan. El grafo produce componentes desconectados y caminos
ausentes de forma automática; instrucciones explícitas amplifican el efecto:
cada artículo de la wiki incluye una sección «Lo que las fuentes no abordan» con
las preguntas adyacentes que el corpus deja abiertas.

\subsubsection*{(c) Anclas conceptuales}

Designar los términos centrales del campo como nodos de alta prioridad
concentra la extracción en lo relevante y reduce el ruido de menciones
incidentales. Para libre competencia chilena, las treinta anclas incluyen:
colusión y bid rigging; abuso de posición dominante; precio predatorio; negativa
de venta y facilidades esenciales; margin squeeze; operación de concentración;
mercado relevante y mercado relevante geográfico; HHI; barreras de entrada;
estándar de prueba; prueba indirecta; argumento rechazado; eficiencias; teoría
del daño; remedio estructural; condición conductual; multa y graduación; nexo
causal; cálculo de daños; before-after y diferencias en diferencias; modelo
contrafactual; DL~211 y su Artículo 3; notificación de concentración.

\subsubsection*{(d) Extracción a nivel de proposición}

Tratar cada sentencia como un único nodo oculta su contenido proposicional. Dos
sentencias alineadas a nivel de documento pueden defender proposiciones
incompatibles a nivel de considerando. El esquema instruye a la IA a extraer,
por cada documento, la lista de proposiciones afirmativas, cada una con tres
campos:

\begin{itemize}
  \item \textbf{Premisa}: el hecho o dato empírico que la sustenta, con
    referencia al considerando.
  \item \textbf{Inferencia}: la conclusión jurídica o económica que el documento
    extrae.
  \item \textbf{Puntero}: referencia exacta al texto.
\end{itemize}

Este nivel de granularidad expone incompatibilidades invisibles a nivel de
documento y constituye el insumo directo del registro de hipótesis.

% =============================================================================
\section{Resultados preliminares}
% =============================================================================

Los hallazgos de esta sección son provisionales. Reflejan el estado del grafo en
la fecha de redacción; el lector puede replicar el pipeline y producir los suyos
sobre el corpus actualizado.

\subsection{Métricas del grafo (Tracks A + C)}

\begin{table}[h!]
\centering
\caption{Métricas del grafo sobre Tracks A + C}
\label{tab:metricas-grafo}
\small
\begin{tabular}{ll}
\toprule
\rowcolor{tabhead}\color{white}\textbf{Métrica} & \color{white}\textbf{Valor} \\
\midrule
Documentos fuente & 520 \\
Nodos & \pend{leer de graphify-out/GRAPH\_REPORT.md} \\
Aristas & \pend{leer de GRAPH\_REPORT.md} \\
Comunidades (Leiden) & \pend{leer de GRAPH\_REPORT.md} \\
Artículos wiki generados & \pend{contar archivos en graphify-out/wiki/} \\
Tiempo de extracción & 15--30 minutos en 1--2 sesiones \\
\bottomrule
\end{tabular}
\end{table}

\noindent
Los valores pendientes se completan a partir de \texttt{graphify-out/GRAPH\_REPORT.md},
que reporta distribución de grado, coeficiente de clustering, longitud media de
camino, entropía entre comunidades y los cinco \emph{god nodes} de mayor grado
con sus vecinos inmediatos.

\subsection{God nodes}

Los \emph{god nodes} son los nodos de mayor grado: los documentos y conceptos a
los que el resto del corpus se conecta con mayor frecuencia. Su centralidad no
la impone un editor humano; emerge del patrón de citas y dependencias
conceptuales en todo el material.\autocite{schrepel2026brain}

\begin{table}[h!]
\centering
\caption{God nodes preliminares del grafo chileno}
\label{tab:god-nodes}
\small
\begin{tabularx}{\linewidth}{>{\raggedright}p{5.5cm} l >{\raggedright\arraybackslash}X}
\toprule
\rowcolor{tabhead}\color{white}\textbf{Nodo} &
  \color{white}\textbf{Tipo} &
  \color{white}\textbf{Tema principal} \\
\midrule
\pend{god node 1} & Sentencia / concepto & Definición de mercado relevante \\
\pend{god node 2} & Sentencia CS & Estándar de prueba en colusión \\
\pend{god node 3} & Sentencia TDLC & Abuso de posición dominante \\
\pend{god node 4} & Concepto & Art.~3 DL~211 \\
\pend{god node 5} & ICG / Resolución TDLC & Instrucción normativa de base \\
\bottomrule
\end{tabularx}
\end{table}

\subsection{Componentes desconectados}

El grafo contiene clusters sin conexiones extraídas hacia el cuerpo principal
del corpus. Son las fronteras del campo y sus callejones doctrinales. Un
componente desconectado puede representar una subcategoría emergente aún no
citada en el mainstream, una línea de argumentación intentada y abandonada, o
un cuerpo normativo que el corpus no ha integrado.\autocite{schrepel2026brain}

\begin{table}[h!]
\centering
\caption{Componentes desconectados preliminares}
\label{tab:desconectados}
\small
\begin{tabularx}{\linewidth}{>{\raggedright}X >{\raggedright\arraybackslash}X}
\toprule
\rowcolor{tabhead}\color{white}\textbf{Cluster} &
  \color{white}\textbf{Hipótesis de aislamiento} \\
\midrule
\pend{cluster aislado 1} &
  Área emergente sin precedente jurisprudencial consolidado \\
\pend{cluster aislado 2} &
  Normativa FNE sin aristas hacia sentencias que la apliquen \\
Acuerdos extrajudiciales pre-2010 &
  Poca práctica de citar AE antiguos como precedente \\
\pend{otro cluster} & \pend{hipótesis} \\
\bottomrule
\end{tabularx}
\end{table}

\subsection{Caminos ausentes}

La ausencia de un camino entre dos nodos que en términos teóricos deberían
estar conectados es un hallazgo: significa que el corpus no ha establecido
explícitamente esa conexión. En el grafo chileno destacan, en estado preliminar:

\begin{itemize}
  \item \textbf{Metodologías econométricas $\leftrightarrow$ sentencias.} Las
    sentencias del TDLC citan o critican modelos econométricos —regresión de
    daños, GUPPI, diferencias en diferencias— pero el grafo muestra que esos
    modelos no tienen nodo propio con aristas entrantes desde las sentencias que
    los evalúan. La «metodología econométrica» existe en el corpus como objeto
    de discusión, no como concepto con masa doctrinal autónoma.
  \item \textbf{Investigaciones FNE $\leftrightarrow$ fallos TDLC} (Track~B
    pendiente). Las sentencias citan a la FNE de origen, pero las
    investigaciones aún no están en el grafo. Incorporado el Track~B, este
    camino debe materializarse; si no lo hace, su ausencia será objeto de
    investigación.
  \item \textbf{Eficiencias $\leftrightarrow$ condenas.} «Eficiencias» aparece
    en informes FNE y en argumentos de defensa, pero el grafo no muestra
    aristas sólidas hacia sentencias donde el argumento haya sido acogido para
    mitigar o excluir la sanción. La ausencia puede reflejar una brecha real o
    un déficit de extracción.
  \item \pend{otro camino ausente detectado en consulta de grafo}.
\end{itemize}

\subsection{Tensiones TDLC/CS}

\begin{table}[h!]
\centering
\caption{Tensiones TDLC/CS identificadas en el grafo (preliminar)}
\label{tab:tensiones}
\small
\begin{tabularx}{\linewidth}{>{\raggedright}p{3.8cm} >{\raggedright}X >{\raggedright\arraybackslash}X}
\toprule
\rowcolor{tabhead}\color{white}\textbf{Punto doctrinal} &
  \color{white}\textbf{Posición TDLC} &
  \color{white}\textbf{Posición CS} \\
\midrule
Prueba indirecta (\emph{plus factors}) en colusión &
  \pend{posición TDLC} & \pend{posición CS} \\
\addlinespace
Definición de mercado relevante en mercados digitales &
  \pend{posición TDLC} & \pend{posición CS} \\
\addlinespace
Factores relevantes en la graduación de la multa &
  \pend{posición TDLC} & \pend{posición CS} \\
\addlinespace
\pend{otra tensión detectada} &
  \pend{posición TDLC} & \pend{posición CS} \\
\bottomrule
\end{tabularx}
\end{table}

\subsection{Silencios doctrinales}

Los silencios son tan informativos como el contenido. En el corpus chileno
emergen, en estado preliminar:

\begin{itemize}
  \item \textbf{Cálculo de daños.} Concepto presente en informes periciales,
    sin masa doctrinal propia en jurisprudencia TDLC. Sugiere ausencia de una
    metodología consolidada sobre cuantificación de sobreprecios.
  \item \textbf{Presunción de inocencia en libre competencia.} Argumento usado
    en defensas, sin pronunciamiento sustantivo del TDLC ni de la CS que lo fije
    para el campo.
  \item \textbf{Responsabilidad civil de terceros.} Pregunta abierta en la
    doctrina académica del corpus; ausente en la jurisprudencia.
  \item \pend{silencio adicional detectado en consulta}.
\end{itemize}

\noindent
El Anexo~\ref{anx:tabla-resumen} consolida estos hallazgos en una tabla única
que conviene para referencia rápida.

% =============================================================================
\section{El protocolo de consulta de seis pasos}
\label{sec:protocolo}
% =============================================================================

El protocolo es una adaptación del propuesto por Schrepel para convertir el
grafo en un instrumento de generación de hipótesis y no solo de recuperación de
información.\autocite{schrepel2026brain} Los seis pasos se ejecutan en
secuencia. Juntos producen una agenda de investigación estructurada: lo que el
campo ha establecido, dónde están sus fronteras y qué conexiones no ha logrado
hacer. Los \emph{prompts} concretos figuran en el Anexo~\ref{anx:consultas}.

\subsubsection*{Paso 1 — Identificar los god nodes}

\begin{codebox}
/graphify query "what are the most connected nodes in the graph"
\end{codebox}

Los god nodes revelan qué trata el campo como fundacional. Un corpus en que
cierta Sentencia TDLC es el nodo más conectado es un corpus organizado en torno
a una teoría particular de la dominancia. Un corpus en que «mercado relevante»
es el nodo más conectado está organizado en torno a la definición de mercado
como eje analítico.

\subsubsection*{Paso 2 — Identificar componentes desconectados}

\begin{codebox}
/graphify query "identify disconnected components"
\end{codebox}

Cada componente desconectado es una pregunta de investigación: un área
emergente, un argumento abandonado, o un cuerpo normativo todavía huérfano de
jurisprudencia que lo aplique.

\subsubsection*{Paso 3 — Consultar caminos ausentes}

\begin{codebox}
/graphify query "identify absent paths in the graph"
\end{codebox}

Tras ubicar god nodes y clusters aislados, consultar por caminos entre pares de
nodos que deberían estar conectados en términos teóricos pero no lo están. La
ausencia es el hallazgo.

\subsubsection*{Paso 4 — Divergencia entre corpora}

Cuando existen dos grafos sobre el mismo dominio, comparar sus god nodes,
comunidades y conexiones. Para el corpus chileno, la comparación natural es
entre el grafo jurisprudencial (Track~A) y el grafo de enforcement FNE
(Track~B): conceptos densos en el primero y aislados en el segundo son
candidatos a investigación sobre brechas entre doctrina judicial y práctica
administrativa.

\subsubsection*{Paso 5 — Verificar conflictos de autoridad}

\begin{codebox}
/graphify query "authority conflicts: TDLC positions reversed by CS"
\end{codebox}

La jerarquía codificada en \texttt{CLAUDE.md} permite que el grafo surfacee
automáticamente estos conflictos. Detectarlos por sistema es más fiable que
confiar en la memoria del investigador.

\subsubsection*{Paso 6 — Chequeo de salud}

\begin{codebox}
/graphify query "concepts mentioned frequently but never defined"
\end{codebox}

Conceptos con muchas aristas entrantes y pocas salientes, sin artículo de wiki
que los defina, son los conocimientos asumidos del campo: términos que todos los
documentos usan pero ninguno explicita. En libre competencia chilena, candidatos
naturales son «eficiencias», «bien jurídico protegido» y «nexo causal».

% =============================================================================
\section{Discusión}
\label{sec:discusion}
% =============================================================================

\subsection{Investigación estructural y práctica diaria}

La separación entre grafo y RAG no es estética: responde a necesidades del
usuario que son genuinamente distintas. El economista-abogado que prepara un
escrito necesita el texto exacto del considerando que respalda un argumento;
para eso, RAG. El mismo profesional, cuando diseña la estrategia de un caso
nuevo, necesita saber cuántas veces el TDLC aceptó prueba indirecta en colusión,
en qué porcentaje la CS revocó esas decisiones y en qué casos el argumento de
eficiencias mitigó la sanción; para eso, el grafo.

Un sistema único que pretenda satisfacer ambos requerimientos —típicamente, un
chatbot con búsqueda semántica— termina sirviendo parcialmente a los dos sin
optimizar a ninguno. La arquitectura propuesta aquí separa instrumentos y les
asigna las preguntas para las que están diseñados.

\subsection{El valor de los silencios}

La sección «Lo que las fuentes no abordan» en cada artículo de wiki es la
expresión visible de una orientación metodológica: en campos estructurados por
debates, las preguntas sin respuesta suelen ser más útiles que las que ya la
tienen.\autocite{schrepel2026brain} La recuperación de contenido —encontrar lo
que el corpus dice sobre un tema— está ampliamente disponible con búsqueda
convencional. Lo que la búsqueda convencional no entrega es el mapa de lo que
el corpus \emph{no} dice. El grafo lo produce como subproducto; consultar por
caminos ausentes, clusters aislados y conceptos no definidos lo convierte en
output primario.

Para el investigador de libre competencia, los silencios tienen un valor
adicional: muestran qué argumentos no tienen respaldo jurisprudencial sólido,
qué metodologías no han sido validadas judicialmente y qué mercados el TDLC y la
FNE nunca han analizado en profundidad. Esa información alimenta tanto la
estrategia de litigio como la agenda de investigación académica.

\subsection{Control de alucinaciones}

El riesgo de alucinación es menor en un sistema basado en grafo que en un
chatbot de propósito general, pero no desaparece. La estructura del grafo
restringe al asistente al contenido ya extraído. La extracción misma, sin
embargo, puede errar: un agente puede inferir una arista inexistente, asignar un
considerando al documento equivocado o crear nodos con etiquetas inconsistentes.

Tres mecanismos mitigan el problema en la implementación descrita:

\begin{enumerate}
  \item \textbf{Diseño del esquema.} Un \texttt{CLAUDE.md} que instruye al
    agente a citar considerandos y referencias exactas reduce las aristas
    inferidas sin ancla textual.
  \item \textbf{Extracción con Gemini.} Usar la API de Gemini como proveedor de
    extracción —en lugar del modelo principal de la sesión— deja libre el
    contexto de Claude Code para la supervisión y la consulta.
  \item \textbf{Verificación manual de hallazgos críticos.} El protocolo de seis
    pasos genera hipótesis; el investigador las verifica en los documentos
    fuente. El grafo detecta; no decide.
\end{enumerate}

\subsection{Limitaciones}

Tres limitaciones conviene señalar de frente. Primero, el grafo no captura
metodologías econométricas con la granularidad que un economista esperaría:
extrae la mención del método, no los supuestos, los datos ni los resultados.
Para análisis metodológico fino, los documentos fuente siguen siendo el lugar
primario de consulta. Segundo, la calidad de extracción sobre documentos
escaneados es desigual: sentencias TDLC anteriores a 2010 en formato imagen
tienen tasas de error más altas. Tercero, el grafo es una fotografía del corpus
en el momento de la extracción; los documentos posteriores no se ven hasta que
se corre \texttt{--update}.

\subsection{Extensiones}

El sistema híbrido más completo integraría Graphify (investigación) y un sistema
RAG vectorial (práctica) sobre el mismo corpus, con una interfaz unificada que
enrute cada pregunta al instrumento adecuado. La incorporación del Track~B
—2.641 documentos de enforcement FNE— es la extensión más urgente: habilita el
Paso~4 del protocolo (divergencia entre corpora) y el rastreo completo de casos
vinculados. Una tercera extensión natural es el despliegue institucional: un
tribunal o una agencia podría publicar una versión consultable de su propio
corpus decisional, permitiendo que cualquier interesado consulte en lenguaje
natural y reciba respuestas trazables a documentos específicos.

% =============================================================================
\section{Conclusión}
% =============================================================================

Este artículo documentó la construcción de un cerebro digital para la
investigación en libre competencia chilena. El pipeline es replicable, de código
abierto y neutro respecto del tamaño del corpus. MarkItDown para conversión,
Graphify para construcción del grafo, y Claude Code como entorno de consulta
son herramientas de licencia MIT que no requieren infraestructura especializada
ni experiencia avanzada en programación. \texttt{CLAUDE.md} es el único punto de
intervención del investigador en la extracción; su diseño cuidadoso es lo que
separa un sistema de recuperación de un instrumento de generación de hipótesis.

Las contribuciones propias son modestas pero específicas. La jerarquía
CS/TDLC/FNE codificada en el esquema refleja particularidades institucionales
que un esquema genérico no captura. La distinción operativa entre grafo
(investigación) y RAG (práctica) responde a una necesidad concreta del
profesional que litiga y que la literatura existente no articula con claridad.
El protocolo de seis pasos sobre el corpus jurisprudencial produce un tipo de
output —agenda de investigación, mapa de silencios, registro de tensiones
TDLC/CS— que la búsqueda convencional no entrega.

El sistema resultante es un primer paso. La incorporación del Track~B, el
despliegue de un sistema RAG, y la depuración de la extracción a nivel de
proposición son tareas que este primer paso vuelve posibles. El repositorio
—con el \texttt{CLAUDE.md} anotado, los scripts de descarga y conversión, y este
artículo— quedará disponible públicamente, para que investigadores de otras
jurisdicciones o áreas del derecho puedan replicar y adaptar el pipeline.

La invitación es a la comunidad: a usar el sistema, a mejorar el esquema, a
identificar los god nodes del propio campo y a interrogar sus silencios.

% =============================================================================
\section*{Atribución y reconocimientos}
\addcontentsline{toc}{section}{Atribución y reconocimientos}
% =============================================================================

La metodología de base proviene de Andrej Karpathy, que describió la práctica
mínima de mantener una carpeta de documentos fuente y usar un asistente de IA
para organizarla y consultarla.\autocite{spisak2026} La habilidad Graphify, que
ejecuta la construcción del grafo, fue desarrollada por Safi
Shamsi.\autocite{shamsi2025graphify} MarkItDown, que convierte PDFs a Markdown,
fue desarrollado por Microsoft.\autocite{microsoft2024markitdown} La guía
metodológica que orienta la arquitectura de investigación de este artículo fue
escrita por Thibault Schrepel.\autocite{schrepel2026brain}

Las contribuciones propias son tres y se enuncian de forma acotada:

\begin{enumerate}
  \item \textbf{Adaptación institucional.} El esquema de Schrepel se traduce al
    sistema chileno: jerarquía CS/TDLC/FNE, tensión tribunal especialista frente
    a generalista, extracción de proposiciones a nivel de considerando.
  \item \textbf{Protocolo de seis pasos.} Adaptado para generar hipótesis sobre
    jurisprudencia de competencia, no solo para recuperar información.
  \item \textbf{Distinción grafo / RAG.} Separación operativa entre instrumento
    de investigación estructural y de práctica diaria, articulada en torno a
    una necesidad concreta del economista-abogado que litiga.
\end{enumerate}

Todo lo demás —la mecánica de Graphify, el registro de hipótesis, el módulo de
diagnóstico— es adaptación para investigación, no trabajo técnico original. La
guía de Schrepel es el lugar donde buscar la fundamentación metodológica completa.

% =============================================================================
% REFERENCIAS
% =============================================================================

\printbibliography[title={Referencias}]

% =============================================================================
% ANEXOS
% =============================================================================

\appendix

\section{Contenido del archivo CLAUDE.md (versión anonimizada)}
\label{anx:claude-md}

El archivo \texttt{CLAUDE.md} es el esquema de diseño que moldea la extracción.
A continuación se reproduce su estructura con las rutas absolutas y la
información personal removidas. Los elementos en \texttt{[...]} son
\emph{placeholders} para que el lector complete con sus propios datos. El cuerpo
del artículo discute las cuatro decisiones de diseño que distinguen este esquema
(Sección~\ref{sec:esquema}).

\begin{codebox}
# Base de Conocimiento — [NOMBRE DEL CAMPO]

## Qué es esto
Una base de conocimiento de [práctica profesional / investigación] sobre
[descripción del dominio]. El corpus cubre [tipos de documentos].
El sistema está orientado a [descripción del usuario].

## Cómo está organizado
- raw/         — documentos fuente en PDF. No modificar.
- raw-md/      — versiones Markdown de raw/. No modificar.
- wiki/        — wiki organizada por temas. La IA la mantiene.
- notes/       — notas personales del investigador. No son fuentes
                 de autoridad y nunca deben citarse como tal.
- outputs/     — minutas, informes y documentos generados.
- code/        — scripts de descarga y transformación.

## Jerarquía de autoridad de las fuentes
1. [Fuente de mayor autoridad] — vinculante
2. [Segunda fuente] — especializado; supeditado a 1
   2b. [Variante normativa de 2] — vinculante para agentes del mercado
   2c. [Acuerdos entre instituciones] — vinculantes para las partes
3. [Autoridad administrativa]
4. [Guías no vinculantes]
5. [Estudios sectoriales]
6. [Informes periciales]
7. [Doctrina académica]
0. Notas personales — nivel cero; sin autoridad jurídica

## Regla de conflicto
Cuando fuentes de diferentes niveles aborden el mismo punto, citar
la de mayor autoridad y señalar cómo confirma, modifica o anula
a la de menor autoridad.

## Ponderación por centralidad de nodo (Node Centrality Weighting)
El grafo rankea los documentos por conexiones. Los más conectados
son "god nodes" y reflejan el consenso del campo. Al generar artículos
de wiki y sintetizar respuestas, dar mayor peso a los god nodes.

## Reglas de la wiki
- Cada tema tiene su propio .md en wiki/
- Estructura fija: (1) Resumen, (2) Jurisprudencia, (3) Enforcement,
  (4) Casos vinculados, (5) Lo que las fuentes no abordan
- Vincular temas con [[nombre-del-tema]]
- Mantener INDEX.md con descripción de una línea por tema
- Al agregar nuevas fuentes, actualizar los artículos relevantes

## Conceptos centrales (anclas de extracción — extraer como nodos
## de alta prioridad)
[Lista de 15-30 conceptos centrales del campo]

## Reglas de extracción de proposiciones
Para cada documento, extraer proposiciones afirmativas con:
- Premisa: hecho empírico con referencia al párrafo
- Inferencia: conclusión jurídica/económica
- Puntero: referencia exacta al texto

## Registro de hipótesis
Cada vez que una consulta revele una contradicción, generar una
hipótesis y registrarla en hypotheses.md con: fecha, consulta de
origen, documentos involucrados, tipo de inconsistencia.

## Módulo de diagnóstico
Después de cada build, calcular propiedades de red del grafo y
reportarlas en GRAPH_REPORT.md:
- Distribución de grado (¿sigue una ley de potencia?)
- Coeficiente de clustering + longitud media de camino
- Entropía entre comunidades
- Traza año a año

## Perspectiva de uso
[Descripción del rol del usuario y sus tres necesidades principales]

## Reglas de respuesta
1. Citas textuales — obligatorio (texto exacto + referencia precisa)
2. Nivel de detalle — intermedio aplicado
3. Sin precedente local — indicarlo explícitamente
4. Ventana temporal — sin corte fijo; último pronunciamiento relevante
5. Precedentes recientes — énfasis en los últimos 5-7 años
\end{codebox}

% ----- Anexo B -----

\section{Ejemplos de consultas al grafo}
\label{anx:consultas}

A continuación se reproducen los \emph{prompts} del protocolo de seis pasos
(Sección~\ref{sec:protocolo}) con espacio para la respuesta. Los marcadores
\pend{respuesta real} deben completarse corriendo el pipeline sobre el corpus
actualizado.

\subsubsection*{Consulta 1 — God nodes}

\begin{codebox}
/graphify query "what are the most connected nodes in the graph"
\end{codebox}

\pend{Salida de Graphify con los cinco nodos de mayor grado y el cluster al
que pertenecen.}

Formato esperado:
\begin{quote}
\itshape
Los cinco nodos de mayor grado son: (1)~[Sentencia TDLC N°X/YYYY, grado~Z];
(2)~[concepto «mercado relevante», grado~W]; (3)~[\dots]; (4)~[\dots];
(5)~[\dots]. El nodo más conectado pertenece al cluster
«[nombre de la comunidad]» e indica que el campo se organiza alrededor de
[interpretación].
\end{quote}

\subsubsection*{Consulta 2 — Componentes desconectados}

\begin{codebox}
/graphify query "identify disconnected components"
\end{codebox}

\pend{Salida de Graphify con los clusters aislados y su tamaño.}

\subsubsection*{Consulta 3 — Caminos ausentes}

\begin{codebox}
/graphify query "identify absent paths between efficiency defenses
  and TDLC condemnation decisions"
\end{codebox}

\pend{Salida de Graphify con los pares de nodos no conectados.}

\subsubsection*{Consulta 4 — Conflictos de autoridad}

\begin{codebox}
/graphify query "authority conflicts: which TDLC positions were
  reversed or modified by the Corte Suprema"
\end{codebox}

\pend{Salida de Graphify con los pares de pronunciamientos contradictorios.}

\subsubsection*{Consulta 5 — Chequeo de salud}

\begin{codebox}
/graphify query "concepts mentioned frequently but never defined:
  show nodes with many incoming edges but no outgoing edges"
\end{codebox}

\pend{Salida de Graphify con los conceptos con masa entrante y sin nodo
de definición.}

% ----- Anexo C -----

\section{Tabla resumen de resultados preliminares}
\label{anx:tabla-resumen}

La Tabla~\ref{tab:resumen-hallazgos} consolida los hallazgos estructurales de
la Sección~5. Se mantiene aquí como referencia rápida; cada fila puede leerse
junto con la sección correspondiente del cuerpo principal.

\begin{longtable}{>{\raggedright}p{3.5cm} >{\raggedright}p{4cm}
    >{\raggedright\arraybackslash}p{5.5cm}}
\caption{Resumen de hallazgos estructurales del grafo (Tracks A + C)}
\label{tab:resumen-hallazgos}\\
\toprule
\rowcolor{tabhead}\color{white}\textbf{Categoría} &
  \color{white}\textbf{Hallazgo} &
  \color{white}\textbf{Interpretación / Hipótesis} \\
\midrule
\endfirsthead
\multicolumn{3}{l}{\small\itshape Continuación de la Tabla~\ref{tab:resumen-hallazgos}}\\
\toprule
\rowcolor{tabhead}\color{white}\textbf{Categoría} &
  \color{white}\textbf{Hallazgo} &
  \color{white}\textbf{Interpretación / Hipótesis} \\
\midrule
\endhead
\bottomrule
\endfoot
God node 1 & \pend{nombre} & Concepto o sentencia fundacional del campo \\
God node 2 & \pend{nombre} & \pend{interpretación} \\
God node 3 & \pend{nombre} & \pend{interpretación} \\
\addlinespace
Componente desconectado 1 & AE pre-2010 & Poca práctica de citar AE como precedente \\
Componente desconectado 2 & \pend{cluster} & \pend{hipótesis} \\
\addlinespace
Camino ausente 1 & Eficiencias $\leftrightarrow$ condenas &
  La jurisprudencia no conecta explícitamente el argumento con el resultado \\
Camino ausente 2 & Metodologías econométricas $\leftrightarrow$ sentencias &
  Los modelos son objeto de discusión, no nodos con masa doctrinal \\
\addlinespace
Silencio 1 & Cálculo de daños &
  Brecha: no hay doctrina consolidada sobre cuantificación \\
Silencio 2 & Presunción de inocencia en libre competencia &
  Argumento mencionado sin pronunciamiento sustantivo \\
\addlinespace
Tensión TDLC/CS 1 & Prueba indirecta en colusión &
  \pend{número de casos; dirección de la tensión} \\
Tensión TDLC/CS 2 & \pend{tensión} & \pend{interpretación} \\
\end{longtable}

\end{document}

% =============================================================================
% RESUMEN DE CAMBIOS RESPECTO DE main.tex
% =============================================================================
%
% ESTRUCTURA
% ----------
% 1. Se eliminó la sección suelta "Nota sobre atribución" entre el abstract y
%    la introducción. Su contenido se reubicó al final, antes de las Referencias,
%    como sección sin numerar "Atribución y reconocimientos". Esto descomprime
%    la entrada al artículo y mejora el flujo Intro → Metodología → Implementación.
%
% 2. Se fusionaron las dos secciones que en el borrador original trataban por
%    separado "CLAUDE.md como instrumento de diseño" e "Implementación en el
%    corpus chileno". En la versión corregida son una sola sección
%    (sec:esquema) con tres bloques: corpus, tensión TDLC/CS, decisiones de
%    diseño del CLAUDE.md. Evita la repetición del corpus y de la jerarquía de
%    autoridad en dos lugares.
%
% 3. La introducción ahora cierra con la subsección sobre Chile (por qué este
%    caso) en lugar de pasar abruptamente a la metodología.
%
% 4. El orden global queda: Introducción → Metodología → Implementación →
%    Resultados → Protocolo de consulta → Discusión → Conclusión → Atribución →
%    Referencias → Anexos. Es el orden pedido.
%
% GRAPHIFY VS RAG
% ---------------
% Se agregó un nuevo entorno `keybox` (tcolorbox) y se construyó un recuadro
% destacado con título "Graphify para investigación vs. RAG vectorial para el
% día a día" al final de la introducción. La tabla comparativa (Tabla 1)
% expandió sus criterios: añade tiempo de consulta y reordena foco/utilidad.
%
% PLACEHOLDERS
% ------------
% Se cargó el paquete `todonotes`. Se define el comando \pend{...}, que renderiza
% un recuadro amarillo discreto con la marca "Pendiente:". Reemplaza todos los
% `[COMPLETAR: ...]` del borrador. La sintaxis permite desactivarlos todos para
% la versión final (cambiar la línea \usepackage{todonotes} a
% \usepackage[disable]{todonotes}).
%
% PROSA
% -----
% Reescritura integral con foco en:
% - Frases más cortas; eliminación de subordinadas innecesarias.
% - Transiciones explícitas entre párrafos.
% - Eliminación de muletillas tipo "El resultado es...", "La característica más
%   relevante es...", "Es importante destacar...".
% - Verbos en voz activa donde antes había nominalizaciones.
% - Variedad léxica: se evitan repeticiones de "el sistema", "el corpus" en
%   párrafos consecutivos.
% - Registro académico aplicado uniforme (sin oscilación divulgativa).
%
% COHESIÓN CON ANEXOS
% -------------------
% Se añadieron referencias cruzadas explícitas entre el cuerpo y los anexos:
%   - Sección de metodología → Anexo A (CLAUDE.md)
%   - Sección de implementación → Anexo A (CLAUDE.md)
%   - Sección de protocolo → Anexo B (consultas)
%   - Sección de resultados → Anexo C (tabla resumen)
% En cada anexo se agregó un párrafo introductorio que lo conecta con el cuerpo.
%
% CORRECCIONES TÉCNICAS
% ---------------------
% - No se alteró ningún comando de Graphify, MarkItDown ni Claude Code.
% - No se modificó la descripción de Gemini API ni del esquema en CLAUDE.md.
% - No se introdujeron afirmaciones nuevas sobre el sistema legal chileno; se
%   reorganizó y comprimió el material existente.
% - Se conservó la bibliografía y todas las \autocite originales.
%
% TIPOGRAFÍA Y FORMATO
% --------------------
% - Nuevo entorno `keybox` (tcolorbox) para el recuadro destacado.
% - Color `accent` añadido para el borde de las notas \pend{}.
% - Atribución sin numerar pero incluida en el índice via \addcontentsline.
% =============================================================================
